\documentclass[11pt,a4paper]{article}
\usepackage[T1]{fontenc}
\usepackage[utf8]{inputenc}
\usepackage{amsmath,amssymb,xcolor}
\usepackage[margin=25mm]{geometry}
\usepackage[colorlinks=true,urlcolor=blue]{hyperref}
\definecolor{nexorablue}{HTML}{2459B5}
\setlength{\parindent}{0pt}
\setlength{\parskip}{7pt}
\setlength{\emergencystretch}{2em}
\renewcommand{\familydefault}{\sfdefault}
\begin{document}

{\small\bfseries\color{nexorablue} NEXORA CPGE}\par
{\small MP / NEX-SLC-03 / Énoncé}\par
{\LARGE\bfseries Un capteur de gain différent de un}\par
\vspace{0.25em}\hrule\vspace{0.6em}

La chaîne directe est $G(p)=6/(1+0{,}40p)$ et le capteur a un gain $b=1/3$. On applique un échelon de consigne unitaire. Calculer le gain et la constante de temps en boucle fermée, l’écart de mesure final et le temps de réponse à $5\,\%$.


\end{document}
